\subsection{Stack Tecnológico}\label{subsec:stack}

La solución Paralegales se estructuró como un ecosistema de tres repositorios independientes, diseñados para operar de forma desacoplada pero coordinada. En conjunto, estos componentes comparten una base funcional apoyada en Firebase para la autenticación y el almacenamiento de datos durante la etapa inicial de la plataforma, mientras que cada capa resuelve un propósito específico dentro de la arquitectura general: los servicios de \textit{backend}, la experiencia pública de usuarios y el portal administrativo.

El repositorio \texttt{paralegales-services} concentra la capa de servicios y automatización del sistema. Está implementado en Go y orientado hacia una arquitectura de microservicios con ejecución \textit{serverless} sobre AWS Lambda. Su diseño incorpora integración con API Gateway para la exposición de \textit{endpoints}, mensajería asíncrona mediante SQS y SNS, e infraestructura como código a través de CloudFormation. Adicionalmente, incluye observabilidad centralizada con CloudWatch y trazabilidad distribuida con AWS X-Ray, lo que facilita el monitoreo operativo de los flujos críticos. Este repositorio integra también servicios externos como Twilio para el envío de mensajes de WhatsApp, así como Firestore y utilidades de Firebase para consultas transaccionales y validación de usuarios. En términos funcionales, soporta procesos como notificaciones, migraciones, análisis de actuaciones judiciales, generación de documentos y otros flujos \textit{event-driven} que dependen de colas de mensajería y funciones desacopladas.

El repositorio \texttt{paralegales-ui} corresponde a la aplicación principal orientada al usuario final. Está desarrollado con Nuxt~3 y Vue~3, combinación que permite una experiencia moderna de \textit{frontend} con capacidades de renderizado del lado del servidor cuando el despliegue lo requiere. Su \textit{stack} incorpora Vuetify como biblioteca de componentes visuales, Pinia para la gestión de estado global, Vee~Validate y Zod para la validación de formularios. En materia de aseguramiento de calidad, el repositorio integra Playwright para pruebas de extremo a extremo (\textit{end-to-end}) y Vitest para pruebas unitarias. Esta capa actúa simultáneamente como \textit{frontend} público y como \textit{Backend-for-Frontend} (BFF) en ciertos flujos, centralizando llamadas al \textit{backend} y gestionando la interacción con Firebase, servicios internos y los componentes de experiencia de usuario.

Por su parte, \texttt{paralegales-admin} implementa el portal administrativo del sistema. Este repositorio está construido con Vue~3, Vite y TypeScript, y utiliza Vuetify para mantener una interfaz coherente y orientada a operaciones internas. Su configuración incorpora Vue Router con generación automática de rutas, Pinia para el estado de la aplicación y Firebase para la autenticación y la consulta de información administrativa. En la práctica, este portal concentra funciones de gestión y supervisión, separadas de la experiencia pública para preservar la claridad operativa y el control de acceso.

En conjunto, estos tres repositorios materializan una arquitectura distribuida en la que el \textit{backend serverless} ejecuta la lógica de negocio y la automatización, mientras que las aplicaciones Vue/Nuxt resuelven los canales de interacción con usuarios finales y administradores. Esta separación de responsabilidades favorece la escalabilidad, el mantenimiento y la evolución independiente de cada capa del sistema. El siguiente apartado profundiza en los patrones arquitectónicos que articulan internamente cada una de estas capas.
